% Fuente: Auxiliar 6, Problema 2
\begingroup
\color{MidnightBlue}
\begin{enumerate}
\item[(a)]
\textbf{Verdadera}, por dualidad fuerte, si ambos problemas tienen soluciones óptimas finitas, entonces sus valores óptimos coinciden: $z^* = w^*$.

\item[(b)]
\textbf{Falsa}. Por dualidad débil, para toda solución factible $x$ de $(P)$ y $y$ de $(D)$, se cumple: $z \leq w$ (en maximización). Por tanto, esta desigualdad está invertida.

\item[(c)]
\textbf{Verdadera}. Es consecuencia directa de la dualidad débil, cualquier solución factible de $(P)$ tiene valor $\leq$ que el óptimo dual, y viceversa.

\item[(d)]
\textbf{Falsa}. Si $(P)$ es infactible, el dual puede ser infactible o no acotado. No hay implicación directa de infactibilidad recíproca.

\item[(e)]
\textbf{Verdadera}. Si $(P)$ tiene solución óptima finita, entonces el dual también tiene una solución óptima finita y los valores óptimos coinciden (teorema de dualidad fuerte).

\item[(f)]
\textbf{Falsa}. Si $(P)$ es no acotado (en maximización), entonces el dual es infactible. No es que sea no acotado, sino que no tiene soluciones factibles.
\end{enumerate}
\endgroup
